Here is a complete LaTeX proof:

```latex
\begin{proof}
Set
\[
r=r(x):=\sqrt{1-x^2}>0,
\qquad
q(x,\psi)=\bigl(x,r\sin\psi,r\cos\psi\bigr),
\]
for $x\in(-1,1)$ and $\psi\in\mathbb{R}/2\pi\mathbb{Z}$.
Also write
\[
W(q):=\frac{d(\mu^++\mu^-)}{dA}(q),
\qquad
w(x,\psi)=W\bigl(q(x,\psi)\bigr).
\]

\begin{enumerate}
\item[(i)] We first compute
\[
q\!\left(x,\psi-\frac{\pi}{2}\right)
=
\left(
x,\,
r\sin\!\left(\psi-\frac{\pi}{2}\right),\,
r\cos\!\left(\psi-\frac{\pi}{2}\right)
\right)
=
\bigl(x,-r\cos\psi,r\sin\psi\bigr).
\]
Since
\[
R_x(x,y,z)=(x,z,-y),
\]
its inverse is
\[
R_x^{-1}(x,y,z)=(x,-z,y).
\]
Therefore
\[
R_x^{-1}\bigl(q(x,\psi)\bigr)
=
R_x^{-1}\bigl(x,r\sin\psi,r\cos\psi\bigr)
=
\bigl(x,-r\cos\psi,r\sin\psi\bigr)
=
q\!\left(x,\psi-\frac{\pi}{2}\right).
\]

Now, because $B=R_x(A)$, the corresponding Green functions satisfy
\[
G_B(q)=G_A\!\bigl(R_x^{-1}(q)\bigr).
\]
Hence
\[
G_B\bigl(q(x,\psi)\bigr)
=
G_A\!\left(R_x^{-1}\bigl(q(x,\psi)\bigr)\right)
=
G_A\!\left(q\!\left(x,\psi-\frac{\pi}{2}\right)\right)
=
F\!\left(x,\psi-\frac{\pi}{2}\right).
\]
This proves
\[
G_B\bigl(q(x,\psi)\bigr)=F\!\left(x,\psi-\frac{\pi}{2}\right).
\]

\item[(ii)] The full cube configuration is invariant under the quarter-turn $R_x$. Hence the total potential
\[
u=2\pi\sum_i G_i
\]
is $R_x$-invariant:
\[
u\circ R_x=u.
\]
Therefore the measures $\mu^\pm$ canonically associated to $u$ are also $R_x$-invariant. (If one views $\mu^\pm$ as the Jordan parts of a signed measure determined by $u$, this follows from uniqueness of the Jordan decomposition.) Since $R_x$ is a rotation, it preserves the round area measure $dA$ as well. Consequently the density
\[
W=\frac{d(\mu^++\mu^-)}{dA}
\]
is $R_x$-invariant.

Next,
\[
R_x\bigl(q(x,\psi)\bigr)
=
R_x\bigl(x,r\sin\psi,r\cos\psi\bigr)
=
\bigl(x,r\cos\psi,-r\sin\psi\bigr).
\]
On the other hand,
\[
q\!\left(x,\psi+\frac{\pi}{2}\right)
=
\left(
x,\,
r\sin\!\left(\psi+\frac{\pi}{2}\right),\,
r\cos\!\left(\psi+\frac{\pi}{2}\right)
\right)
=
\bigl(x,r\cos\psi,-r\sin\psi\bigr).
\]
Thus
\[
R_x\bigl(q(x,\psi)\bigr)=q\!\left(x,\psi+\frac{\pi}{2}\right).
\]
Using the $R_x$-invariance of $W$, we obtain
\[
w\!\left(x,\psi+\frac{\pi}{2}\right)
=
W\!\left(q\!\left(x,\psi+\frac{\pi}{2}\right)\right)
=
W\!\left(R_x\bigl(q(x,\psi)\bigr)\right)
=
W\bigl(q(x,\psi)\bigr)
=
w(x,\psi).
\]
Hence
\[
w\!\left(x,\psi+\frac{\pi}{2}\right)=w(x,\psi)
\]
for all $x$ and $\psi$.

\item[(iii)] Differentiate $q(x,\psi)=\bigl(x,r\sin\psi,r\cos\psi\bigr)$, using
\[
r'(x)=\frac{d}{dx}\sqrt{1-x^2}=-\frac{x}{\sqrt{1-x^2}}=-\frac{x}{r}.
\]
Then
\[
\partial_x q
=
\left(
1,\,
-\frac{x\sin\psi}{r},\,
-\frac{x\cos\psi}{r}
\right),
\qquad
\partial_\psi q
=
\bigl(0,r\cos\psi,-r\sin\psi\bigr).
\]
Their cross product is
\[
\partial_x q\times \partial_\psi q
=
\begin{vmatrix}
\mathbf{i} & \mathbf{j} & \mathbf{k}\\[2pt]
1 & -\dfrac{x\sin\psi}{r} & -\dfrac{x\cos\psi}{r}\\[6pt]
0 & r\cos\psi & -r\sin\psi
\end{vmatrix}.
\]
Computing the determinant gives
\[
\partial_x q\times \partial_\psi q
=
\left(
x\sin^2\psi+x\cos^2\psi,\,
r\sin\psi,\,
r\cos\psi
\right)
=
\bigl(x,r\sin\psi,r\cos\psi\bigr)
=
q(x,\psi).
\]
Since $q(x,\psi)\in S^2$, we have
\[
\bigl|\partial_x q\times \partial_\psi q\bigr|=|q(x,\psi)|=1.
\]
Therefore the surface area element is
\[
dA=\bigl|\partial_x q\times \partial_\psi q\bigr|\,dx\,d\psi=dx\,d\psi.
\]
So the Jacobian is $1$.

\item[(iv)] To avoid confusing the parameter $x$ with an ambient coordinate derivative, let us denote the ambient Euclidean coordinates by $(X,Y,Z)$. The vector field corresponding to $\partial_\phi$ on $S^2$ is
\[
\partial_\phi=-Y\,\partial_X+X\,\partial_Y.
\]
At the point
\[
q(x,\psi)=\bigl(X,Y,Z\bigr)=\bigl(x,r\sin\psi,r\cos\psi\bigr),
\]
this tangent vector is
\[
(-Y,X,0)=(-r\sin\psi,x,0).
\]

Now the pushed-forward coordinate vector fields on the parameter domain are exactly
\[
\partial_x q
=
\left(
1,\,
-\frac{x\sin\psi}{r},\,
-\frac{x\cos\psi}{r}
\right),
\qquad
\partial_\psi q
=
\bigl(0,r\cos\psi,-r\sin\psi\bigr).
\]
We claim that
\[
(-r\sin\psi,x,0)
=
-r\sin\psi\,\partial_x q
+
\frac{x\cos\psi}{r}\,\partial_\psi q.
\]
Indeed,
\[
-r\sin\psi\,\partial_x q
+
\frac{x\cos\psi}{r}\,\partial_\psi q
\]
\[
=
-r\sin\psi
\left(
1,\,
-\frac{x\sin\psi}{r},\,
-\frac{x\cos\psi}{r}
\right)
+
\frac{x\cos\psi}{r}\,(0,r\cos\psi,-r\sin\psi)
\]
\[
=
\left(
-r\sin\psi,\,
x\sin^2\psi+x\cos^2\psi,\,
x\sin\psi\cos\psi-x\sin\psi\cos\psi
\right)
=
(-r\sin\psi,x,0).
\]
Thus, as an operator on functions of $(x,\psi)$,
\[
\partial_\phi
=
-r\sin\psi\,\partial_x+\frac{x\cos\psi}{r}\,\partial_\psi.
\]
Substituting back $r=\sqrt{1-x^2}$ yields
\[
\partial_\phi
=
-\sqrt{1-x^2}\,\sin\psi\,\partial_x
+
\frac{x\cos\psi}{\sqrt{1-x^2}}\,\partial_\psi.
\]

This proves all four claims.
\end{enumerate}
\end{proof}
```

If you want, I can also reformat this into a more polished theorem-proof style matching the rest of your document.